% SVY 5430 Biyomedikal Verilerin Analizi ve Yapay Zeka — Dönem Projesi Rapor Şablonu
% Derleme: pdflatex main && bibtex main && pdflatex main && pdflatex main   (ya da Overleaf'e yükleyin)
\documentclass[11pt,a4paper]{article}

\usepackage[utf8]{inputenc}
\usepackage[T1]{fontenc}
\usepackage[turkish,shorthands=:!]{babel}
\usepackage{lmodern}
\usepackage[margin=2.5cm]{geometry}
\usepackage{graphicx}
\usepackage{booktabs}
\usepackage{amsmath}
\usepackage{hyperref}
\usepackage{xcolor}
\usepackage{caption}
\usepackage{float}
\hypersetup{colorlinks=true, linkcolor=black, citecolor=blue!60!black, urlcolor=blue!60!black}
\graphicspath{{sekiller/}}

% ---- Künye: burayı doldurun ------------------------------------------------
\newcommand{\projekodu}{BVYZ-XX}
\newcommand{\projebaslik}{Proje Başlığı: Hangi Veriyle Neyi Tahmin Ediyorsunuz}
\newcommand{\ogrenci}{Ad Soyad}
\newcommand{\ogrencino}{Öğrenci No}
\newcommand{\depo}{https://github.com/kullanici/depo}
% ----------------------------------------------------------------------------

\title{\projebaslik\\[0.4em]{\large SVY 5430 Biyomedikal Verilerin Analizi ve Yapay Zeka --- Dönem Projesi (\projekodu)}}
\author{\ogrenci\ (\ogrencino)\\Fırat Üniversitesi, Yazılım ve Bilişim Araştırma Enstitüsü\\\small Kod ve defterler: \url{\depo}}
\date{\today}

\begin{document}
\maketitle

\begin{abstract}
En fazla 200 kelime. Klinik soru, veri (kaynak, örnek sayısı), yöntem (ön işleme, model, doğrulama), ana sonuç (en iyi modelin ROC-AUC ve kalibrasyonu) ve açıklanabilirlik bulgusunun bir cümlelik özeti. Son cümle: sınırlılık ve klinik anlamı.
\end{abstract}

\section{Giriş ve Klinik Soru}
Problemi bir klinisyene anlatır gibi yazın: Hangi hastalar, hangi sonuç, neden önemli? Ardından soruyu makine öğrenmesi diline çevirin (sınıflandırma / regresyon / kümeleme / segmentasyon) ve etiketin nasıl tanımlandığını açıkça belirtin. Literatürden en az üç çalışmaya atıf yapın~\cite{steyerberg2010}.

\section{Veri}
\subsection{Kaynak ve Erişim}
Veri setinin adı, kaynağı, lisansı ve nasıl indirildiği. Hasta verisi için etik/erişim koşulları.

\subsection{Tanımlayıcı Özet}
Örnek sayısı, öznitelik sayısı, hedef değişken ve sınıf oranı. Tablo~\ref{tab:veri} şablonunu kullanın.

\begin{table}[H]
\centering
\caption{Veri setinin özeti.}
\label{tab:veri}
\begin{tabular}{lrr}
\toprule
Özellik & Değer & Not \\
\midrule
Örnek (hasta / kayıt) & \ldots & \\
Öznitelik sayısı & \ldots & \\
Pozitif sınıf oranı & \ldots\,\% & \\
Eksik değer oranı (en yüksek sütun) & \ldots\,\% & \\
\bottomrule
\end{tabular}
\end{table}

\subsection{Keşifçi Analiz}
En fazla iki şekil: hedefe göre en ayırt edici değişkenler ve (varsa) PCA/UMAP görünümü. Şekil~\ref{fig:eda} gibi.

\begin{figure}[H]
\centering
% \includegraphics[width=0.8\textwidth]{eda.png}
\fbox{\parbox{0.8\textwidth}{\centering\vspace{2cm} sekiller/eda.png \vspace{2cm}}}
\caption{Keşifçi analiz: \ldots}
\label{fig:eda}
\end{figure}

\section{Yöntem}
\subsection{Ön İşleme}
Eksik veri mekanizması ve imputasyon kararı, ölçekleme, kategorik kodlama, dengesiz sınıf stratejisi. Her kararın \emph{neden} verildiğini bir cümleyle yazın.

\subsection{Doğrulama Tasarımı}
Bölme stratejisi (hasta-bazlı mı, zaman-farkında mı), katlama sayısı, sızıntıya karşı alınan önlemler (Pipeline). Bu bölüm kısa ama zorunludur.

\subsection{Modeller}
Temel model (lojistik regresyon) ve en az bir ileri model. Hiperparametre arama yöntemi. Kullandığınız kütüphaneler ve sürümleri.

\subsection{Değerlendirme Ölçütleri}
ROC-AUC, PR-AUC, duyarlılık/özgüllük (seçilen eşikte), kalibrasyon (eğri ve Brier). Eşiği neden o değerde seçtiğinizi klinik maliyetle gerekçelendirin.

\section{Sonuçlar}
\begin{table}[H]
\centering
\caption{Model başarımı (5 katlı çapraz doğrulama, ortalama $\pm$ s.s.).}
\label{tab:sonuc}
\begin{tabular}{lcccc}
\toprule
Model & ROC-AUC & PR-AUC & Brier & Duyarlılık@eşik \\
\midrule
Lojistik regresyon & \ldots & \ldots & \ldots & \ldots \\
\ldots & \ldots & \ldots & \ldots & \ldots \\
\bottomrule
\end{tabular}
\end{table}

ROC ve kalibrasyon eğrilerini tek şekilde yan yana verin. Alt-grup analizi (cinsiyet, yaş grubu) yaptıysanız burada raporlayın.

\section{Açıklanabilirlik}
En az bir küresel (permütasyon önemi, SHAP özet) ve bir yerel (tek hasta için SHAP şelale / Grad-CAM) açıklama~\cite{lundberg2017shap}. Açıklamanın klinik olarak anlamlı olup olmadığını tartışın: model doğru yere mi bakıyor?

\section{Tartışma}
Bulguların klinik anlamı, literatürle karşılaştırma, sınırlılıklar (veri boyutu, tek merkez, etiket kalitesi), dağıtım kayması riski ve bu modelin kliniğe girmesi için ne gerektiği~\cite{collins2024tripod}.

\section{Sonuç}
Üç–dört cümle. Ne öğrendiniz, bir sonraki adım ne olurdu?

\section*{Tekrarlanabilirlik Beyanı}
Kod ve defterler: \url{\depo}. Rastgelelik tohumu, ortam dosyası ve veri indirme talimatı depoda yer alır. Veri dosyaları depoya konmamıştır.

\bibliographystyle{plain}
\bibliography{kaynaklar}

\end{document}
